% Generated by geometry_pipeline.latex_renderer. Do not edit by hand.

\chapter{Кути. Суміжні та вертикальні кути}

\begin{tcolorbox}[colback=boxbglight,colframe=primaryblue,coltitle=white,fonttitle=\bfseries\sffamily,title={СТАТУС АВТОМАТИЧНОЇ ЗБІРКИ},arc=3mm,boxrule=1.2pt,breakable]

Цей файл згенеровано зі структурованих блоків. Позначені місця рисунків очікують векторної реконструкції та редакторського перегляду.

\end{tcolorbox}

% block: geo.p013.s13.definition-angle; printed source pages: 13

\paragraph{13. Кут.}

Фігура, утворена двома променями ($OA$ і $OB$, рис. 8), що виходять з однієї точки, називається \textbf{кутом}. Промені, які утворюють кут, називаються \textbf{сторонами кута}, а точка, з якої вони виходять, — \textbf{вершиною кута}. Сторони кута слід уявляти собі нескінченно продовженими від вершини.

Кут зазвичай позначають трьома буквами, з яких середня ставиться біля вершини, а крайні — біля будь-яких точок на сторонах; наприклад, кажуть: «кут $AOB$» або «кут $BOA$» (рис. 8). Проте можна позначати кут і однією буквою, поставленою біля вершини, якщо при цій вершині немає інших кутів. Іноді ми позначатимемо кут цифрою, поставленою всередині кута біля вершини.

Частина площини, обмежена сторонами кута, розглядається як така, що лежить \textbf{всередині кута}; решта площини лежить \textbf{поза кутом}; перша становить \textbf{внутрішню область кута}, друга — \textbf{зовнішню область} його.

Якщо з вершини кута (рис. 8) проведемо всередині нього будь-які промені $OD,\allowbreak OE \dots$, то утворені при цьому кути $AOD,\allowbreak DOE,\allowbreak EOB \dots$ розглядаються як частини кута $AOB$.

Слово «кут» на письмі часто замінюють символом $\angle$.

\begin{center}
% Extracted from the reviewed manual pilot; edit the source reference first.
\begin{tikzpicture}[scale=1.1, >=Stealth]
    \coordinate (O) at (0,0);
    \coordinate (A) at (4,2.5);
    \coordinate (B) at (4.5,0);
    \coordinate (D) at (3.8,1.6);
    \coordinate (E) at (4.2,0.8);

    \draw[->, very thick, color=primaryblue] (O) -- (A) node[above right] {$A$};
    \draw[->, very thick, color=primaryblue] (O) -- (B) node[right] {$B$};

    \draw[->, thick, dashed, color=gray] (O) -- (D) node[right] {$D$};
    \draw[->, thick, dashed, color=gray] (O) -- (E) node[right] {$E$};

    \filldraw[black] (O) circle (2pt) node[left] {$O$};
\end{tikzpicture}

{\small\textbf{Рис. 8. Кут AOB і внутрішні промені OD та OE}}\\[6pt]
\end{center}

\begin{feynmanbox}[Кут — це міра повороту, а не просто лінії на папері]
Коли ми дивимося на креслення в зошиті, кут здається застиглим трикутним шматочком площини. Проте в реальному фізичному світі кут — це \textbf{динамічний процес обертання}.
Уяви стрілку годинника, циркуль чи промінь обертового маяка. Спочатку промінь був спрямований уздовж рейки $OB$. Потім його повернули навколо шарніра $O$ до положення $OA$.
Величина кута показує, \textbf{наскільки сильно повернувся промінь}. Самі промені нескінченні, тому їхня довжина на папері ніяк не впливає на кут!
\end{feynmanbox}

\begin{mediabox}[Побачити через рух: кут як поворот]
Повертай промінь навколо вершини й спостерігай, як положення променя визначає величину кута:
\href{https://radchr.github.io/Algebra/geometry-angle-rotation.html}{https://radchr.github.io/Algebra/geometry-angle-rotation.html}
\end{mediabox}

% block: geo.p013.s14.angle-equality; printed source pages: 13

\paragraph{14. Рівність і нерівність кутів.}

Два кути вважаються \textbf{рівними}, якщо при накладанні вони можуть суміститися. Припустимо, наприклад, що ми накладаємо кут $AOB$ на кут $A_1O_1B_1$ (рис. 9) так, щоб вершина $O$ збіглася з $O_1$, сторона $OB$ пішла по $O_1B_1$ і щоб кути накрили один одного своїми внутрішніми областями. Якщо при цьому сторона $OA$ суміститься з $O_1A_1$, то кути рівні; якщо ж сторона $OA$ пройде всередині кута $A_1O_1B_1$ або поза ним, то кути не рівні, до того ж меншим буде той з них, який становить частину іншого кута.

\begin{center}
% Extracted from the reviewed manual pilot; edit the source reference first.
\begin{tikzpicture}[scale=0.95, >=Stealth]
    % Перший кут AOB
    \coordinate (O) at (0,0);
    \coordinate (A) at (2.4,2.1);
    \coordinate (B) at (3.2,0);
    \draw[->, very thick, color=primaryblue] (O) -- (A) node[above left] {$A$};
    \draw[->, very thick, color=primaryblue] (O) -- (B) node[right] {$B$};
    \filldraw (O) circle (2pt) node[left] {$O$};

    % Другий кут A1O1B1 з накладанням
    \coordinate (O1) at (5.5,0);
    \coordinate (A1) at (7.9,2.1);
    \coordinate (B1) at (8.7,0);
    \coordinate (A_in) at (7.6,1.4);
    \draw[->, very thick, color=primaryblue] (O1) -- (A1) node[above right] {$A_1$};
    \draw[->, very thick, color=primaryblue] (O1) -- (B1) node[right] {$B_1$};
    \draw[->, thick, dashed, color=warningred] (O1) -- (A_in) node[right] {\small $OA$ (всередині)};
    \filldraw (O1) circle (2pt) node[below left] {$O_1$};

\end{tikzpicture}

{\small\textbf{Рис. 9. Порівняння кутів способом накладання}}\\[6pt]
\end{center}

\motionlink{Рівні кути можна сумістити}%
{Перенеси й поверни копію кута до точного суміщення з еталоном та простеж, що його величина не змінюється.}%
{https://radchr.github.io/Algebra/geometry-equal-angles-motion.html}%
{qr_equal_angles_motion.png}

% block: geo.p014.s15.angle-sum; printed source pages: 14

\paragraph{15. Сума кутів.}

\textbf{Сумою кутів} $AOB$ і $A_1O_1B_1$ (рис. 10) називається такий кут, який отримують у такий спосіб. Будуємо кут $MNP$, що дорівнює першому даному куту $AOB$, і до нього приєднуємо кут $PNQ$, що дорівнює другому даному куту $A_1O_1B_1$, так, щоб в обох кутів виявилася спільна вершина $N$ і спільна сторона $NP$, а внутрішні області кутів були розташовані по різні боки від цієї спільної сторони $NP$. Тоді кут $MNQ$ і буде сумою кутів $AOB$ та $A_1O_1B_1$. Подібним чином може бути складена сума трьох і більше кутів.

Сума кутів, як і сума відрізків прямої, має \textbf{переставну} та \textbf{сполучну} властивості.

З поняття про суму кутів виводяться поняття про їхню \textbf{різницю}, \textbf{добуток} і \textbf{частку}.

Часто доводиться говорити про такий промінь, який ділить даний кут навпіл; такому променю дали спеціальну назву: \textbf{бісектриса} (рис. 11).

\begin{center}
% Extracted from the reviewed manual pilot; edit the source reference first.
\begin{tikzpicture}[scale=0.95, >=Stealth]
    % Кут 1
    \draw[->, thick] (0,0) node[left] {$O$} -- (2,1.5) node[above] {$B$};
    \draw[->, thick] (0,0) -- (2.3,0) node[right] {$A$};
    \draw (0.6,0) arc (0:37:0.6);

    % Кут 2
    \draw[->, thick] (3.5,0) node[left] {$O_1$} -- (5.5,1.0) node[above] {$B_1$};
    \draw[->, thick] (3.5,0) -- (5.8,0) node[right] {$A_1$};
    \draw (4.1,0) arc (0:26:0.6);

    % Сума
    \coordinate (N) at (7.5,0);
    \draw[->, very thick, color=primaryblue] (N) node[below left] {$N$} -- (10.2,0) node[right] {$M$};
    \draw[->, thick, dashed] (N) -- (9.8,1.6) node[above right] {$P$};
    \draw[->, very thick, color=primaryblue] (N) -- (9.0,2.4) node[above] {$Q$};
    \draw[color=primaryblue] (8.2,0) arc (0:58:0.7);

\end{tikzpicture}

{\small\textbf{Рис. 10. Додавання двох кутів}}\\[6pt]
% Extracted from the reviewed manual pilot; edit the source reference first.
\begin{tikzpicture}[scale=1.0, >=Stealth]
    \coordinate (O) at (0,0);
    \coordinate (A) at (4,2.4);
    \coordinate (B) at (4.5,0);
    \coordinate (L) at (4.4,1.2);

    \draw[->, thick] (O) -- (A) node[above right] {$A$};
    \draw[->, thick] (O) -- (B) node[right] {$B$};
    \draw[->, very thick, color=feynmangreen] (O) -- (L) node[right] {\textbf{Бісектриса $OL$}};

    \draw[color=feynmangreen] (0.9,0) arc (0:15.5:0.9);
    \draw[color=feynmangreen] (0.9,0.25) arc (15.5:31:0.9);
    \filldraw[black] (O) circle (2pt) node[left] {$O$};
\end{tikzpicture}

{\small\textbf{Рис. 11. Бісектриса кута}}\\[6pt]
\end{center}

% block: geo.p014.s16.extended-angles; printed source pages: 14, 15

\paragraph{16. Розширення поняття про кут.}

Під час знаходження суми кутів можуть трапитися деякі особливі випадки, які корисно розглянути окремо:

\begin{enumerate}
\item Може статися, що після додавання кількох кутів, наприклад трьох: $AOB,\allowbreak BOC$ і $COD$ (рис. 12), сторона $OD$ кута $COD$ стане продовженням сторони $OA$ кута $AOB$. Ми отримаємо фігуру, утворену двома променями ($OA$ і $OD$), що виходять з однієї точки ($O$) і є доповняльними променями (становлять продовження один одного). Таку фігуру також прийнято називати кутом (\textbf{розгорнутим}). Внутрішня область розгорнутого кута становить половину площини.
\item Може статися, що після додавання кількох кутів, наприклад п'яти: $AOB,\allowbreak BOC,\allowbreak COD,\allowbreak DOE$ і $EOA$ (рис. 13), сторона $OA$ останнього кута суміститься зі стороною $OA$ першого кута. Фігура, утворена такими суміщеними променями (разом з усією площиною навколо їхньої спільної вершини $O$), також називається кутом (\textbf{повним}).
\item Нарешті, будуючи суму кутів, ми можемо не лише заповнити всю площину навколо спільної вершини, а й перекрити її вдруге, втретє тощо. Така сума кутів дорівнює одному повному куту, доданому до деякого кута, двом повним кутам із деяким кутом тощо.
\end{enumerate}

\begin{center}
% Extracted from the reviewed manual pilot; edit the source reference first.
\begin{tikzpicture}[scale=1.0, >=Stealth]
        \coordinate (O) at (0,0);
        \coordinate (A) at (3,0);
        \coordinate (B) at (1.5,1.8);
        \coordinate (C) at (-1.2,1.8);
        \coordinate (D) at (-3,0);

        \draw[->, thick] (O) -- (A) node[right] {$A$};
        \draw[->, thick] (O) -- (B) node[above right] {$B$};
        \draw[->, thick] (O) -- (C) node[above left] {$C$};
        \draw[->, thick] (O) -- (D) node[left] {$D$};

        \draw[color=primaryblue, thick] (0.8,0) arc (0:180:0.8);
        \filldraw (O) circle (2pt) node[below] {$O$};
    \end{tikzpicture}

{\small\textbf{Рис. 12. Розгорнутий кут як сума частин}}\\[6pt]
% Extracted from the reviewed manual pilot; edit the source reference first.
\begin{tikzpicture}[scale=0.9, >=Stealth]
        \coordinate (O) at (0,0);
        \draw[->, thick] (O) -- (2.5,0) node[right] {$A$};
        \draw[->, thick] (O) -- (1.5,2.0) node[above right] {$B$};
        \draw[->, thick] (O) -- (-1.2,2.0) node[above left] {$C$};
        \draw[->, thick] (O) -- (-2.0,-0.8) node[left] {$D$};
        \draw[->, thick] (O) -- (0.5,-2.0) node[below] {$E$};

        \draw[color=primaryblue, thick, ->] (0.6,0) arc (0:350:0.6);
        \filldraw (O) circle (2pt) node[below left=2pt] {$O$};
    \end{tikzpicture}

{\small\textbf{Рис. 13. Повний кут навколо точки O}}\\[6pt]
\end{center}

\begin{warningbox}[Розгорнутий кут — це просто пряма лінія?]
Учні часто кажуть: «Розгорнутий кут — це просто пряма». Це груба помилка!
Пряма лінія не має виділеної точки (вершини). А розгорнутий кут має \textbf{фіксовану вершину} $O$ і \textbf{два протилежні промені}, що виходять із неї. Без фіксованої вершини кут не існує!
\end{warningbox}

% block: geo.p015.s17.central-angle; printed source pages: 15

\paragraph{17. Центральний кут і дві його властивості.}

Кут ($AOB$, рис. 14), утворений двома радіусами, називається \textbf{центральним кутом}; про такий кут і дугу, укладену між його сторонами, кажуть, що вони відповідають один одному.

Центральні кути відносно відповідних їм дуг мають такі дві властивості:
\emph{В одному колі або в рівних колах:}
\begin{enumerate}
\item \textbf{Якщо центральні кути рівні, то й відповідні їм дуги рівні.}
\item \textbf{Обернено: якщо дуги рівні, то й відповідні їм центральні кути рівні.}
\end{enumerate}

\emph{Доведення.} Нехай $\angle AOB = \angle COD$ (рис. 15); потрібно довести, що дуги $AB$ і $CD$ також рівні.
Уявімо, що сектор $AOB$ ми повернули навколо центра $O$ так, щоб радіус $OA$ сумістився з $OC$. Тоді внаслідок рівності кутів радіус $OB$ суміститься з $OD$; отже, сумістяться й дуги $AB$ та $CD$, тобто вони рівні. Друга властивість доводиться аналогічно способом накладання.

\begin{center}
% Extracted from the reviewed manual pilot; edit the source reference first.
\begin{tikzpicture}[scale=0.9, >=Stealth]
    % Рис. 14
    \draw[thick] (0,0) circle (1.5);
    \coordinate (O) at (0,0);
    \coordinate (A) at (-1.1,-1.0);
    \coordinate (B) at (1.1,-1.0);
    \draw[thick] (O) -- (A) node[below left] {$A$};
    \draw[thick] (O) -- (B) node[below right] {$B$};
    \filldraw (O) circle (1.5pt) node[above] {$O$};

    % Рис. 15
    \begin{scope}[xshift=6.5cm]
        \draw[thick] (0,0) circle (1.5);
        \coordinate (O2) at (0,0);
        \coordinate (A2) at (-0.6,1.4);
        \coordinate (B2) at (1.2,0.9);
        \coordinate (C2) at (1.4,-0.4);
        \coordinate (D2) at (0.2,-1.5);
        \draw[thick] (O2) -- (A2) node[above left] {$A$};
        \draw[thick] (O2) -- (B2) node[above right] {$B$};
        \draw[thick] (O2) -- (C2) node[right] {$C$};
        \draw[thick] (O2) -- (D2) node[below right] {$D$};
        \draw[->, thick, color=mediapurple] (1.7,0.7) arc (20:-15:1.7);
        \filldraw (O2) circle (1.5pt) node[left] {$O$};
    \end{scope}
\end{tikzpicture}

{\small\textbf{Рис. 14. Центральний кут}}\\[6pt]
{\small\textbf{Рис. 15. Рівні центральні кути й відповідні дуги}}\\[6pt]
\end{center}

\begin{mediabox}[Побачити через рух: центральний кут і дуга]
Змінюй два рівні центральні кути в рівних колах і спостерігай, чому відповідні їм дуги також рівні:
\href{https://radchr.github.io/Algebra/geometry-central-angle-arc.html}{https://radchr.github.io/Algebra/geometry-central-angle-arc.html}
\end{mediabox}

% block: geo.p015.s18.angle-degrees; printed source pages: 15, 16

\paragraph{18. Градуси: дуговий і кутовий.}

Уявімо, що деяке коло поділено на 360 рівних частин і з кожної точки поділу проведено радіуси до центра. Тоді навколо центра утвориться 360 маленьких центральних кутів, які, відповідаючи рівним дугам, мають бути рівними між собою. Кожна з утворених дуг називається \textbf{дуговим градусом}, а кожен із відповідних кутів при центрі — \textbf{кутовим градусом}.

Отже, дуговий градус є $\frac{1}{360}$ частиною кола (або $\frac{1}{180}$ частиною півкола), а кутовий градус — це центральний кут, що відповідає дуговому градусу. Оскільки 360 таких кутів утворюють повний кут, то кутовий градус є $\frac{1}{360}$ частиною повного кута (або $\frac{1}{180}$ частиною розгорнутого кута).

Градуси поділяються на 60 рівних частин — \textbf{мінут} (позначаються $'$), а мінути — на 60 \textbf{секунд} ($''$).

% block: geo.p016.s19.central-angle-arc; printed source pages: 16

\paragraph{19. Відповідність між центральними кутами й дугами.}

Нехай $AOB$ — деякий кут (рис. 16). Опишемо між його сторонами з вершини $O$ довільним радіусом дугу $CD$. Тоді кут $AOB$ буде центральним кутом, що відповідає дузі $CD$. Припустимо, що ця дуга містить 20 дугових градусів. Тоді й кут $AOB$ поділиться на 20 кутових градусів. Загалом: \textbf{кут вимірюється відповідною йому дугою}. Наприклад:
\[
\angle AOB = 20^\circ 10' 15''.
\]

\begin{center}
% Extracted from the reviewed manual pilot; edit the source reference first.
\begin{tikzpicture}[scale=1.1, >=Stealth]
    \coordinate (O) at (0,0);
    \draw[thick] (O) -- (4,2.2) node[right] {$A$};
    \draw[thick] (O) -- (4.2,0) node[right] {$B$};
    \draw[very thick, color=primaryblue] (2.8,0) arc (0:29:2.8) node[midway, right=2pt] {\small $CD$};
    \node at (2.9, 1.6) {$C$};
    \node at (2.9, -0.3) {$D$};

    % Радіальні поділки
    \foreach \a in {3, 6, 9, 12, 15, 18, 21, 24, 27} {
        \draw[thin, color=gray] (O) -- (\a:2.8);
    }
    \filldraw (O) circle (2pt) node[left] {$O$};
\end{tikzpicture}

{\small\textbf{Рис. 16. Вимірювання кута відповідною дугою}}\\[6pt]
\end{center}

% block: geo.p016.s20.protractor; printed source pages: 16

\paragraph{20. Транспортир.}

Прилад для вимірювання кутів (рис. 17) являє собою півкруг, шкала якого поділена на 180 градусів. Щоб виміряти кут $DCE$, центр півкруга суміщають із вершиною кута, а базовий промінь приладу — зі стороною $CE$. Позначка на шкалі навпроти другої сторони покаже величину кута. За допомогою транспортира можна також побудувати кут заданої градусної міри, наприклад $90^\circ$, $45^\circ$ або $30^\circ$.

\begin{center}
% Extracted from the reviewed manual pilot; edit the source reference first.
\begin{tikzpicture}[scale=1.0, >=Stealth]
    % Корпус транспортира
    \draw[thick, fill=primaryblue!6] (-3.2,0) -- (3.2,0) arc (0:180:3.2) -- cycle;
    \draw[thick, fill=white] (-2.2,0) -- (2.2,0) arc (0:180:2.2) -- cycle;
    \filldraw (0,0) circle (2pt) node[below] {$C$};

    % Градусні поділки
    \foreach \a in {0,10,...,180} {
        \draw (\a:2.8) -- (\a:3.2);
    }
    \foreach \a in {0,30,...,180} {
        \draw[thick] (\a:2.6) -- (\a:3.2);
        \node[font=\tiny] at (\a:2.4) {\a$^\circ$};
    }

    % Вимірюваний кут DCE = 50°
    \draw[very thick, color=primaryblue, ->] (0,0) -- (3.5,0) node[right] {$E$};
    \draw[very thick, color=warningred, ->] (0,0) -- (50:3.5) node[above right] {$D$};
    \draw[color=warningred, thick] (1.1,0) arc (0:50:1.1);
    \node[color=warningred, font=\small\bfseries] at (25:1.45) {$50^\circ$};

\end{tikzpicture}

{\small\textbf{Рис. 17. Вимірювання кута транспортиром}}\\[6pt]
\end{center}

% block: geo.p016.s21.angle-types; printed source pages: 16, 17

\paragraph{21. Прямий, гострий і тупий кути.}

Кут у $90^\circ$ (половина розгорнутого кута) називають \textbf{прямим кутом}; кут, менший від прямого, називають \textbf{гострим}, а більший за прямий (але менший від розгорнутого) — \textbf{тупим} (рис. 18).

Усі прямі кути рівні між собою. Величину прямого кута в класичних підручниках іноді позначають буквою $d$ (від фр. \emph{droit} — прямий).

\begin{center}
% Extracted from the reviewed manual pilot; edit the source reference first.
\begin{tikzpicture}[scale=0.9, >=Stealth]
    % Прямий кут
    \begin{scope}[xshift=0cm]
        \draw[thick] (0,0) -- (2.2,0);
        \draw[thick] (0,0) -- (0,2.2);
        \draw (0.4,0) -- (0.4,0.4) -- (0,0.4);
        \draw[dashed] (1.5,0) arc (0:90:1.5);
        \node at (0.9,0.9) {$90^\circ$};
        \node at (1.0, -0.5) {\small\textbf{прямий}};
    \end{scope}

    % Гострий кут
    \begin{scope}[xshift=4.5cm]
        \draw[thick] (0,0) -- (2.2,0);
        \draw[thick] (0,0) -- (1.5,2.0);
        \draw[color=primaryblue, thick] (0.8,0) arc (0:53:0.8);
        \node at (1.0, -0.5) {\small\textbf{гострий}};
    \end{scope}

    % Тупий кут
    \begin{scope}[xshift=9.0cm]
        \draw[thick] (0,0) -- (2.2,0);
        \draw[thick] (0,0) -- (-1.2,1.8);
        \draw[color=warningred, thick] (0.8,0) arc (0:124:0.8);
        \node at (0.5, -0.5) {\small\textbf{тупий}};
    \end{scope}

\end{tikzpicture}

{\small\textbf{Рис. 18. Прямий, гострий і тупий кути}}\\[6pt]
\end{center}

% block: geo.p017.s22.adjacent-angles; printed source pages: 17

\paragraph{22. Суміжні кути та їхні властивості.}

Два кути ($AOB$ і $BOC$, рис. 19) називаються \textbf{суміжними}, якщо одна сторона в них спільна, а дві інші сторони є доповняльними променями (становлять продовження одна одної).

Оскільки такі кути в сумі утворюють розгорнутий кут, то:
\[
\mathbf{Сума\ двох\ суміжних\ кутів\ дорівнює\ 180^\circ} \quad (\text{або } 2d).
\]

Для будь-якого кута можна побудувати два суміжні з ним кути (рис. 20). Кути, суміжні з одним і тим самим кутом, рівні між собою.

Якщо один із суміжних кутів прямий ($90^\circ$, рис. 21), то й другий дорівнює $180^\circ - 90^\circ = 90^\circ$.
\emph{Наслідок:} Якщо при перетині двох прямих один із чотирьох кутів є прямим, то й інші три кути — прямі.

\begin{center}
% Extracted from the reviewed manual pilot; edit the source reference first.
\begin{tikzpicture}[scale=1.1, >=Stealth]
    \coordinate (O) at (0,0);
    \coordinate (A) at (-3,0);
    \coordinate (C) at (3,0);
    \coordinate (B) at (1.5,2.2);

    \draw[thick] (A) -- (C);
    \draw[->, very thick, color=primaryblue] (O) -- (B) node[above right] {$B$};

    \draw[color=primaryblue, thick] (0.7,0) arc (0:56:0.7) node[midway, right=2pt] {$\beta$};
    \draw[color=warningred, thick] (-0.7,0) arc (180:56:0.7) node[midway, left=2pt] {$\alpha$};

    \filldraw (A) circle (1.5pt) node[below] {$A$};
    \filldraw (O) circle (2pt) node[below] {$O$};
    \filldraw (C) circle (1.5pt) node[below] {$C$};

\end{tikzpicture}

{\small\textbf{Рис. 19. Суміжні кути}}\\[6pt]
% Extracted from the reviewed manual pilot; edit the source reference first.
\begin{tikzpicture}[scale=1.0, >=Stealth]
    \coordinate (O) at (0,0);
    \coordinate (A) at (3,0);
    \coordinate (C) at (-2.5,0);
    \coordinate (B) at (1.8,2.2);
    \coordinate (D) at (-1.8,-2.2);

    \draw[thick] (C) node[left] {$C$} -- (A) node[right] {$A$};
    \draw[thick] (D) node[left] {$D$} -- (B) node[right] {$B$};
    \filldraw (O) circle (2pt) node[below right=2pt] {$O$};

    \node at (1.2, 0.4) {$\angle AOB$};
    \node at (-0.8, 0.5) {$\angle BOC$};
    \node at (0.3, -0.8) {$\angle AOD$};

\end{tikzpicture}

{\small\textbf{Рис. 20. Два кути, суміжні з одним кутом}}\\[6pt]
% Extracted from the reviewed manual pilot; edit the source reference first.
\begin{tikzpicture}[scale=0.9, >=Stealth]
    \draw[thick, <->] (-2.5,0) -- (2.5,0) node[right] {$A$};
    \draw[thick, <->] (0,-2.5) -- (0,2.5) node[above] {$B$};
    \node at (-2.7, 0) {$C$};
    \node at (0, -2.7) {$D$};
    \filldraw (0,0) circle (2pt) node[below left] {$O$};

    \draw (0.4,0) -- (0.4,0.4) -- (0,0.4);
    \draw (-0.4,0) -- (-0.4,0.4) -- (0,0.4);
    \draw (0.4,0) -- (0.4,-0.4) -- (0,-0.4);
    \draw (-0.4,0) -- (-0.4,-0.4) -- (0,-0.4);

    \node at (0.8,0.8) {$90^\circ$};
    \node at (-0.8,0.8) {$90^\circ$};
    \node at (-0.8,-0.8) {$90^\circ$};
    \node at (0.8,-0.8) {$90^\circ$};
\end{tikzpicture}

{\small\textbf{Рис. 21. Перпендикулярні прямі}}\\[6pt]
\end{center}

\begin{scaffoldbox}[Властивість суміжних кутів]
\begin{enumerate}
\item Сторони $OA$ і $OC$ утворюють пряму лінію, тобто розгорнутий кут $\angle AOC$.
\item Градусна міра розгорнутого кута за домовленістю становить $180^\circ$.
\item Внутрішній промінь $OB$ ділить розгорнутий кут на дві частини: $\angle AOB$ та $\angle BOC$.
\item Оскільки ціле дорівнює сумі частин: $\angle AOB + \angle BOC = \angle AOC = 180^\circ$.
\end{enumerate}
\end{scaffoldbox}

\begin{warningbox}[Будь-які кути з сумою $180^\circ$ є суміжними?]
\textbf{Ні в якому разі!} Якщо один кут $70^\circ$ накреслено на початку зошита, а кут $110^\circ$ — на іншій сторінці, їхня сума справді $180^\circ$, але вони \textbf{не суміжні}.
Суміжність обов'язково вимагає спільної вершини та спільної сторони!
\end{warningbox}

\begin{mediabox}[Побачити через рух: суміжні кути]
Рухай спільний промінь і перевіряй, що сума двох суміжних кутів залишається $180^\circ$:
\href{https://radchr.github.io/Algebra/geometry-adjacent-angles.html}{https://radchr.github.io/Algebra/geometry-adjacent-angles.html}
\end{mediabox}

% block: geo.p018.s23.perpendicular-oblique; printed source pages: 18

\paragraph{23. Перпендикуляр і похила.}

Спільна сторона ($OB$) двох суміжних кутів називається \textbf{похилою} до прямої ($AC$), на якій лежать дві інші сторони, коли суміжні кути не рівні між собою (рис. 22, ліворуч). Якщо ж суміжні кути рівні (рис. 22, праворуч), тобто є прямими, ця спільна сторона називається \textbf{перпендикуляром} до прямої. Точка $O$ називається відповідно \textbf{основою похилої} або \textbf{основою перпендикуляра}.

Дві прямі, що перетинаються під прямим кутом, називаються \textbf{перпендикулярними}:
\[
AC \perp BD.
\]

\begin{center}
% Extracted from the reviewed manual pilot; edit the source reference first.
\begin{tikzpicture}[scale=0.9, >=Stealth]
    \draw[thick, <->] (-2.5,0) -- (2.5,0) node[right] {$A$};
    \draw[thick, <->] (0,-2.5) -- (0,2.5) node[above] {$B$};
    \node at (-2.7, 0) {$C$};
    \node at (0, -2.7) {$D$};
    \filldraw (0,0) circle (2pt) node[below left] {$O$};

    \draw (0.4,0) -- (0.4,0.4) -- (0,0.4);
    \draw (-0.4,0) -- (-0.4,0.4) -- (0,0.4);
    \draw (0.4,0) -- (0.4,-0.4) -- (0,-0.4);
    \draw (-0.4,0) -- (-0.4,-0.4) -- (0,-0.4);

    \node at (0.8,0.8) {$90^\circ$};
    \node at (-0.8,0.8) {$90^\circ$};
    \node at (-0.8,-0.8) {$90^\circ$};
    \node at (0.8,-0.8) {$90^\circ$};
\end{tikzpicture}

{\small\textbf{Рис. 21. Перпендикулярні прямі}}\\[6pt]
% Extracted from the reviewed manual pilot; edit the source reference first.
\begin{tikzpicture}[scale=0.95, >=Stealth]
    % Похила
    \draw[thick] (-2.5,0) node[below] {$A$} -- (2.0,0) node[below] {$C$};
    \draw[very thick, color=primaryblue, ->] (0,0) -- (1.5,2.0) node[above right] {$B$};
    \filldraw (0,0) circle (2pt) node[below] {$O$};
    \node at (0, -0.9) {\small\textbf{Похила $OB$}};

    % Перпендикуляр
    \begin{scope}[xshift=6cm]
        \draw[thick] (-2.0,0) node[below] {$A$} -- (2.0,0) node[below] {$C$};
        \draw[very thick, color=feynmangreen, ->] (0,0) -- (0,2.0) node[above] {$B$};
        \draw (0.3,0) -- (0.3,0.3) -- (0,0.3);
        \filldraw (0,0) circle (2pt) node[below] {$O$};
        \node at (0, -0.9) {\small\textbf{Перпендикуляр $OB$}};
    \end{scope}

\end{tikzpicture}

{\small\textbf{Рис. 22. Похила і перпендикуляр до прямої}}\\[6pt]
\end{center}

% block: geo.p018.s24.set-square; printed source pages: 18

\paragraph{24. Креслярський косинець.}

Для побудови перпендикуляра застосовують креслярський косинець із прямим кутом (рис. 23). Його прикладають катетом до лінійки, суміщеної з прямою $AB$, і пересувають до потрібної точки $C$ на прямій або $D$ поза нею.

\emph{Основна властивість:} Через будь-яку точку прямої можна провести перпендикуляр до цієї прямої, і до того ж тільки один.

\begin{center}
% Extracted from the reviewed manual pilot; edit the source reference first.
\begin{tikzpicture}[scale=0.9, >=Stealth]
    % Лінійка
    \draw[fill=gray!20, thick] (-3,-0.5) rectangle (3.5,0);
    \draw[thick] (-3,0) -- (3.5,0);
    \node at (-2.7, 0.3) {$A$};
    \node at (3.2, 0.3) {$B$};

    % Косинець
    \draw[fill=primaryblue!15, very thick] (0,0) -- (0,3.0) -- (2.2,0) -- cycle;
    \draw (0,0.4) -- (0.4,0.4) -- (0.4,0);
    \filldraw (0,0) circle (2pt) node[below] {$C$};
    \filldraw (0,1.8) circle (2pt) node[left] {$D$};
    \node at (0, 3.3) {$E$};
\end{tikzpicture}

{\small\textbf{Рис. 23. Побудова перпендикуляра косинцем}}\\[6pt]
\end{center}

% block: geo.p018.s25.vertical-angles; printed source pages: 18, 19

\paragraph{25. Вертикальні кути.}

Два кути називаються \textbf{вертикальними}, якщо сторони одного кута є доповняльними променями сторін другого. При перетині двох прямих $AB$ і $CD$ (рис. 24) утворюються дві пари вертикальних кутів ($AOD$ і $COB$, $AOC$ і $DOB$).

\[
\mathbf{Вертикальні\ кути\ рівні\ між\ собою:}\quad \angle AOD = \angle BOC.
\]

\begin{center}
% Extracted from the reviewed manual pilot; edit the source reference first.
\begin{tikzpicture}[scale=1.1, >=Stealth]
    \coordinate (O) at (0,0);
    \coordinate (A) at (-2.5,-1.2);
    \coordinate (B) at (2.5,1.2);
    \coordinate (C) at (-2,1.5);
    \coordinate (D) at (2,-1.5);

    \draw[thick, <->] (A) node[below left] {$A$} -- (B) node[above right] {$B$};
    \draw[thick, <->] (C) node[above left] {$C$} -- (D) node[below right] {$D$};
    \filldraw (O) circle (2pt) node[above=2pt] {$O$};

    \node at (-1.5, 0.2) {$\angle 1$};
    \node at (1.5, -0.2) {$\angle 2$};
    \node at (0, 1.0) {$\angle 3$};
    \node at (0, -1.0) {$\angle 4$};
\end{tikzpicture}

{\small\textbf{Рис. 24. Вертикальні кути при перетині двох прямих}}\\[6pt]
\end{center}

\begin{feynmanbox}[Модель ножиць]
Візьми канцелярські ножиці. Коли розкриваються леза (кут 1), ручки ножиць позаду гвинтика (кут 2) розходяться на такий самий кут!
Чому? Бо лезо і ручка — це єдина суцільна пряма лінія. Обидві лінії перетинаються в точці $O$. Якщо один кут збільшується на $10^\circ$, його суміжний сусід зменшується на $10^\circ$, що змушує вертикальний кут збільшитися на ті самі $10^\circ$!
\end{feynmanbox}

\begin{scaffoldbox}[Властивість вертикальних кутів]
Нехай при перетині прямих $AB$ і $CD$ утворилися вертикальні кути $\angle AOD$ та $\angle BOC$. Обидва ці кути є суміжними з одним і тим самим кутом $\angle DOB$.
\begin{enumerate}
\item $\angle AOD + \angle DOB = 180^\circ \implies \angle AOD = 180^\circ - \angle DOB$.
\item $\angle BOC + \angle DOB = 180^\circ \implies \angle BOC = 180^\circ - \angle DOB$.
\item Праві частини однакові, тому й ліві частини рівні: $\angle AOD = \angle BOC$. Теорему доведено!
\end{enumerate}
\end{scaffoldbox}

\begin{mediabox}[Побачити через рух: вертикальні кути]
Поверни копію кута навколо точки перетину на $180^\circ$ і побач, як вона суміщається з вертикальним кутом:
\href{https://radchr.github.io/Algebra/geometry-vertical-half-turn.html}{https://radchr.github.io/Algebra/geometry-vertical-half-turn.html}
\end{mediabox}

% block: geo.p019.s26.common-vertex-angles; printed source pages: 19

\paragraph{26. Кути зі спільною вершиною.}

\begin{enumerate}
\item \textbf{Сума кутів, які мають спільну вершину й розташовані по один бік від прямої} (рис. 25), дорівнює $180^\circ$ (утворюють розгорнутий кут):
\end{enumerate}
   \[
\angle 1 + \angle 2 + \angle 3 + \angle 4 = 180^\circ.
\]
\begin{enumerate}
\item \textbf{Сума кутів навколо спільної вершини} (рис. 26) дорівнює $360^\circ$ (утворюють повний кут):
\end{enumerate}
   \[
\angle 1 + \angle 2 + \angle 3 + \angle 4 + \angle 5 = 360^\circ.
\]
\begin{enumerate}
\item \textbf{Ознака суміжних кутів:} Якщо два кути мають спільну сторону, а їхня сума дорівнює $180^\circ$, то дві інші сторони утворюють пряму лінію (тобто ці кути є суміжними, рис. 27).
\end{enumerate}

\begin{center}
% Extracted from the reviewed manual pilot; edit the source reference first.
\begin{tikzpicture}[scale=0.9, >=Stealth]
        \coordinate (O) at (0,0);
        \draw[thick] (-3,0) node[left] {$A$} -- (3,0) node[right] {$E$};
        \draw[thick, ->] (O) -- (-2,1.8) node[above left] {$B$};
        \draw[thick, ->] (O) -- (0,2.2) node[above] {$C$};
        \draw[thick, ->] (O) -- (2,1.8) node[above right] {$D$};
        \filldraw (O) circle (2pt) node[below] {$O$};
    \end{tikzpicture}

{\small\textbf{Рис. 25. Кути по один бік від прямої}}\\[6pt]
% Extracted from the reviewed manual pilot; edit the source reference first.
\begin{tikzpicture}[scale=0.85, >=Stealth]
        \coordinate (O) at (0,0);
        \draw[thick, ->] (O) -- (2.2,0) node[right] {$A$};
        \draw[thick, ->] (O) -- (1.2,2.0) node[above right] {$B$};
        \draw[thick, ->] (O) -- (-1.5,1.8) node[above left] {$C$};
        \draw[thick, ->] (O) -- (-2.2,-0.5) node[left] {$D$};
        \draw[thick, ->] (O) -- (0.3,-2.0) node[below] {$E$};
        \filldraw (O) circle (2pt) node[below left=2pt] {$O$};
    \end{tikzpicture}

{\small\textbf{Рис. 26. Кути навколо спільної вершини}}\\[6pt]
% Extracted from the reviewed manual pilot; edit the source reference first.
\begin{tikzpicture}[scale=1.0, >=Stealth]
        \coordinate (O) at (0,0);
        \coordinate (A) at (-3,0);
        \coordinate (C) at (3,0);
        \coordinate (B) at (1.2,2.2);

        \draw[thick] (A) node[below] {$A$} -- (C) node[below] {$C$};
        \draw[very thick, color=primaryblue, ->] (O) -- (B) node[above right] {$B$};
        \filldraw (O) circle (2pt) node[below] {$O$};
    \end{tikzpicture}

{\small\textbf{Рис. 27. Ознака суміжних кутів}}\\[6pt]
\end{center}

\clearpage

\section{Вправи}

% block: geo.p019.ex01.supplementary-angle; printed source pages: 19

\begin{problem}[Вправа 1]

Один із суміжних кутів дорівнює $38^\circ 29'$. Знайдіть величину другого кута.

\end{problem}

\begin{scaffoldbox}[Розв'язання або доведення]
\emph{Розв'язання:}
Сума двох суміжних кутів дорівнює $180^\circ$. Представимо $180^\circ$ як $179^\circ 60'$:
\[
\text{Шуканий кут} = 179^\circ 60' - 38^\circ 29' = (179 - 38)^\circ (60 - 29)' = 141^\circ 31'.
\]
\emph{Відповідь:} $141^\circ 31'$.
\end{scaffoldbox}

% block: geo.p019.ex02.straight-line-test; printed source pages: 19

\begin{problem}[Вправа 2]

Два кути $\angle ABC = 100^\circ 20'$ і $\angle CBD = 79^\circ 40'$ мають спільну вершину $B$, спільну сторону $BC$ і не перекриваються. Чи утворюють сторони $AB$ і $BD$ пряму лінію?

\end{problem}

\begin{scaffoldbox}[Розв'язання або доведення]
\emph{Розв'язання:}
Знайдемо суму цих кутів:
\[
\angle ABC + \angle CBD = 100^\circ 20' + 79^\circ 40' = 179^\circ 60' = 180^\circ.
\]
Оскільки сума двох кутів зі спільною стороною дорівнює $180^\circ$, за властивістю § 26.3 їхні сторони $AB$ і $BD$ утворюють розгорнутий кут, тобто лежать на одній прямій лінії.
\emph{Відповідь:} Так, утворюють пряму лінію.
\end{scaffoldbox}

% block: geo.p019.ex03.construct-bisector; printed source pages: 19

\begin{problem}[Вправа 3]

Побудуйте довільний кут і за допомогою транспортира та лінійки проведіть його бісектрису.

\end{problem}

\begin{scaffoldbox}[Розв'язання або доведення]
\emph{Розв'язання:}
\begin{enumerate}
\item Накреслити довільний кут $\angle AOB$.
\item Прикласти транспортир центром до вершини $O$ і виміряти градусну міру кута $\alpha$.
\item Обчислити положення бісектриси: $\frac{\alpha}{2}$.
\item Поставити точку на шкалі транспортира навпроти значення $\frac{\alpha}{2}$ і за допомогою лінійки сполучити вершину $O$ з цією точкою променем. Цей промінь і є бісектрисою.
\end{enumerate}
\end{scaffoldbox}

% block: geo.p020.ex04.adjacent-bisectors; printed source pages: 20

\begin{problem}[Вправа 4]

Доведіть, що бісектриси двох суміжних кутів взаємно перпендикулярні.

\end{problem}

\begin{scaffoldbox}[Розв'язання або доведення]
\emph{Доведення:}
Нехай $\angle AOB = \alpha$ і $\angle BOC = \beta$ — два суміжні кути. Тоді:
\[
\alpha + \beta = 180^\circ.
\]
Проведемо їхні бісектриси $OK$ та $OL$. За означенням бісектриси:
\[
\angle KOB = \frac{\alpha}{2}, \quad \angle BOL = \frac{\beta}{2}.
\]
Кут між бісектрисами дорівнює сумі цих частин:
\[
\angle KOL = \angle KOB + \angle BOL = \frac{\alpha}{2} + \frac{\beta}{2} = \frac{\alpha + \beta}{2} = \frac{180^\circ}{2} = 90^\circ.
\]
Оскільки кут між бісектрисами дорівнює $90^\circ$, вони взаємно перпендикулярні: $OK \perp OL$. Доведено!
\end{scaffoldbox}

\begin{mediabox}[Побачити через рух: бісектриси суміжних кутів]
Змінюй суміжні кути й перевіряй, що кут між їхніми бісектрисами незмінно дорівнює $90^\circ$:
\href{https://radchr.github.io/Algebra/geometry-adjacent-bisectors.html}{https://radchr.github.io/Algebra/geometry-adjacent-bisectors.html}
\end{mediabox}

% block: geo.p020.ex05.vertical-bisectors; printed source pages: 20

\begin{problem}[Вправа 5]

Доведіть, що бісектриси двох вертикальних кутів лежать на одній прямій (є доповняльними променями).

\end{problem}

\begin{scaffoldbox}[Розв'язання або доведення]
\emph{Доведення:}
Нехай при перетині двох прямих утворилися вертикальні кути $\angle 1$ та $\angle 2$, причому $\angle 1 = \angle 2 = \alpha$. Суміжний з ними кут позначимо $\angle 3$, тоді $\angle 3 = 180^\circ - \alpha$.
Проведемо бісектриси обох вертикальних кутів. Кожна бісектриса ділить свій кут на кути по $\frac{\alpha}{2}$.
Розглянемо кут між двома променями бісектрис, який складається з половини першого кута, кута $\angle 3$ і половини другого кута:
\[
\text{Кут між променями} = \frac{\alpha}{2} + \angle 3 + \frac{\alpha}{2} = \frac{\alpha}{2} + (180^\circ - \alpha) + \frac{\alpha}{2} = 180^\circ.
\]
Оскільки кут між променями бісектрис дорівнює $180^\circ$, ці промені утворюють розгорнутий кут, тобто лежать на одній прямій лінії. Доведено!
\end{scaffoldbox}

\begin{mediabox}[Побачити через рух: бісектриси вертикальних кутів]
Рухай одну з прямих і спостерігай, як бісектриси вертикальних кутів залишаються протилежними променями:
\href{https://radchr.github.io/Algebra/geometry-vertical-bisectors.html}{https://radchr.github.io/Algebra/geometry-vertical-bisectors.html}
\end{mediabox}

% block: geo.p020.ex06.opposite-rays; printed source pages: 20

\begin{problem}[Вправа 6]

Якщо при точці $O$ прямої $AB$ (рис. 24) побудувати по різні боки від цієї прямої рівні кути $\angle AOD = \angle BOC$, то промені $OD$ і $OC$ утворюють одну пряму лінію.

\end{problem}

\begin{center}
% Extracted from the reviewed manual pilot; edit the source reference first.
\begin{tikzpicture}[scale=1.1, >=Stealth]
    \coordinate (O) at (0,0);
    \coordinate (A) at (-2.5,-1.2);
    \coordinate (B) at (2.5,1.2);
    \coordinate (C) at (-2,1.5);
    \coordinate (D) at (2,-1.5);

    \draw[thick, <->] (A) node[below left] {$A$} -- (B) node[above right] {$B$};
    \draw[thick, <->] (C) node[above left] {$C$} -- (D) node[below right] {$D$};
    \filldraw (O) circle (2pt) node[above=2pt] {$O$};

    \node at (-1.5, 0.2) {$\angle 1$};
    \node at (1.5, -0.2) {$\angle 2$};
    \node at (0, 1.0) {$\angle 3$};
    \node at (0, -1.0) {$\angle 4$};
\end{tikzpicture}

{\small\textbf{Рис. 24. Вертикальні кути при перетині двох прямих}}\\[6pt]
\end{center}

\begin{scaffoldbox}[Розв'язання або доведення]
\emph{Доведення:}
Кут $\angle AOC$ і кут $\angle BOC$ є суміжними, отже:
\[
\angle AOC + \angle BOC = 180^\circ.
\]
За умовою $\angle BOC = \angle AOD$. Підставимо $\angle AOD$ замість $\angle BOC$:
\[
\angle AOC + \angle AOD = 180^\circ.
\]
Кути $\angle AOC$ та $\angle AOD$ мають спільну вершину $O$, спільну сторону $OA$ і лежать по різні боки від неї, а їхня сума дорівнює $180^\circ$. Отже, за § 26.3 промені $OD$ і $OC$ утворюють розгорнутий кут, тобто лежать на одній прямій лінії. Доведено!
\end{scaffoldbox}

% block: geo.p020.ex07.two-intersecting-lines; printed source pages: 20

\begin{problem}[Вправа 7]

На рис. 24 з точки $O$ проведено промені $OA$, $OD$, $OB$, $OC$. Доведіть, що коли $\angle AOC = \angle DOB$ та $\angle AOD = \angle COB$, то промені $OA$ і $OB$, а також $OD$ і $OC$ є доповняльними променями (утворюють дві прямі лінії, що перетинаються в точці $O$).

\end{problem}

\begin{center}
% Extracted from the reviewed manual pilot; edit the source reference first.
\begin{tikzpicture}[scale=1.1, >=Stealth]
    \coordinate (O) at (0,0);
    \coordinate (A) at (-2.5,-1.2);
    \coordinate (B) at (2.5,1.2);
    \coordinate (C) at (-2,1.5);
    \coordinate (D) at (2,-1.5);

    \draw[thick, <->] (A) node[below left] {$A$} -- (B) node[above right] {$B$};
    \draw[thick, <->] (C) node[above left] {$C$} -- (D) node[below right] {$D$};
    \filldraw (O) circle (2pt) node[above=2pt] {$O$};

    \node at (-1.5, 0.2) {$\angle 1$};
    \node at (1.5, -0.2) {$\angle 2$};
    \node at (0, 1.0) {$\angle 3$};
    \node at (0, -1.0) {$\angle 4$};
\end{tikzpicture}

{\small\textbf{Рис. 24. Вертикальні кути при перетині двох прямих}}\\[6pt]
\end{center}

\begin{scaffoldbox}[Розв'язання або доведення]
\emph{Доведення:}
Усі чотири кути мають спільну вершину $O$ і заповнюють усю площину навколо точки $O$. Їхня загальна сума дорівнює повному куту ($360^\circ$):
\[
\angle AOC + \angle DOB + \angle AOD + \angle COB = 360^\circ.
\]
Враховуючи рівності умов $\angle AOC = \angle DOB$ та $\angle AOD = \angle COB$, маємо:
\[
2 \cdot \angle AOC + 2 \cdot \angle AOD = 360^\circ \implies 2(\angle AOC + \angle AOD) = 360^\circ \implies \angle AOC + \angle AOD = 180^\circ.
\]
Оскільки сума кутів $\angle AOC$ та $\angle AOD$ зі спільною стороною $OA$ дорівнює $180^\circ$, сторони $OC$ і $OD$ є доповняльними променями і утворюють пряму лінію $CD$.
Аналогічно, $\angle AOC + \angle COB = 180^\circ$, тому $OA$ і $OB$ є доповняльними променями та утворюють пряму лінію $AB$.
Отже, обидві пари променів утворюють дві прямі $AB$ і $CD$, що перетинаються в точці $O$. Доведено!
\end{scaffoldbox}
